\booksection{附：安卓手機記憶體拓展解析}

\chapterintro{買手機時，你常能看到“12GB+8GB運存”這樣的宣傳——好像憑空多了8GB執行記憶體。但它真的能讓手機變快嗎？會不會反而傷手機？這一節我們就把“記憶體拓展”這項技術徹底講清楚。}

\booksubsection{一、本質定義}

安卓手機的“記憶體拓展”（又稱記憶體擴充、記憶體融合、RAM Plus等）\textbf{並非真正物理增加執行記憶體晶片}，而是一項系統級軟體技術：透過將手機部分內部儲存空間（ROM/UFS快閃記憶體）動態劃轉，模擬為虛擬執行記憶體（RAM）使用。

打個比方：物理RAM是“高速工作臺”，拓展記憶體則是“臨時儲物間”。當工作臺不夠用時，系統會把暫時不用的資料搬到儲物間，需要時再搬回來。儲物間容量大，但來回搬運肯定比直接放在工作臺上慢。

\booksubsection{二、技術原理}

主流廠商的實現方式主要有兩種：

\subsubsection*{zRAM壓縮}

在RAM內部劃分一塊壓縮資料空間，將不活躍的應用程式碼進行壓縮儲存，壓縮率約2:1。資料仍在RAM中，速度快、不損耗快閃記憶體，但需CPU參與壓縮解壓。

\subsubsection*{swap分割區}

直接在UFS快閃記憶體中劃分固定空間作為虛擬記憶體，將後台程式資料寫入儲存。依賴儲存頻寬，讀寫速度遠低於實體記憶體。

廠商宣傳的“12GB+8GB運存”通常是12GB物理RAM+8GB快閃記憶體虛擬記憶體的組合。值得注意的是，zRAM預設並不佔用機身儲存，只有在zRAM被資料完全填滿時，系統才會在儲存上建立交換檔案，且這種情況極少發生。

\booksubsection{三、速度差距是關鍵瓶頸}

這項技術的根本侷限在於速度差異。目前主流的UFS 4.0快閃記憶體順序讀速約4200MB/s，僅為LPDDR5X實體記憶體的1/10–1/15。因此在輕負荷的zRAM壓縮階段速度接近真實RAM，但一旦進入高負載場景，速度會驟降到幾十分之一。

\booksubsection{四、適用場景與收益邊界}

記憶體拓展並非萬能，效果因機型和使用場景而異：

\begin{itemize}
  \item \textbf{有用的情況}：6GB實體記憶體+UFS 3.1以上、剩餘空間$\geq$30GB的機型，開啟3GB拓展後後臺保活數可多3-4個，過載率降低68\%，適合外賣小哥、網約車司機等重度多工使用者。
  \item \textbf{日常場景}：對微信多開、網頁標籤保活、即時通訊訊息響應等改善顯著，後臺應用冷啟動耗時明顯縮短。
  \item \textbf{有限場景}：對大型遊戲載入、4K影片剪輯等高頻寬場景提升有限。
  \item \textbf{有害的情況}：8GB記憶體+UFS 2.2且剩餘空間不足的機型，開啟後遊戲幀率波動反而擴大，刷短影片多耗電，切換卡頓。
  \item \textbf{無效的情況}：16GB以上大記憶體機型收益幾乎為零，最多多保活1個APP，卻會讓CPU負載增加12\%、發熱上升2-3℃。對RAM大於12GB的機型，建議直接關閉以避免無意義的壓縮計算。
\end{itemize}

\booksubsection{五、各廠商命名與開啟方式}

不同品牌叫法各異：小米叫“記憶體擴充”、華為稱“記憶體融合”、三星叫“RAM Plus”、vivo命名為“Extended RAM”、一加稱作“Memory Extension”。該功能屬系統原生能力，與安卓底層版本無直接繫結關係。

一般開啟路徑為：進入【設定】【更多設定】找到【記憶體】或【效能】子項開啟“記憶體拓展”並選擇容量檔位（通常2GB–7GB），重啟後生效。開啟後系統會顯示如“8GB+4GB”的記憶體組合。

\begin{tcbwarn}{注意事項}
\begin{itemize}
  \item 記憶體拓展會佔用部分機身儲存空間，請確保剩餘儲存不少於20-30GB，否則可能影響系統流暢度。
  \item 快閃記憶體作為虛擬記憶體會有讀寫損耗，雖然現代UFS快閃記憶體壽命較長，但長期高頻swap仍可能加速儲存老化。
  \item 如果手機本身實體記憶體較小（如4GB及以下），開啟記憶體拓展只能緩解燃眉之急，無法從根本上解決卡頓問題。
\end{itemize}

\end{tcbwarn}

\begin{tcbinfo}{小結}
記憶體拓展是一項成熟的系統級最佳化技術，能在硬體不變的前提下提升中低端機型的多工能力與後臺保活率，但它受限於快閃記憶體與實體記憶體的巨大速度鴻溝，並非“萬能神藥”。是否開啟、開多大，需結合自身機型的實體記憶體大小、快閃記憶體規格、剩餘儲存空間及使用習慣來判斷。

\end{tcbinfo}

\begin{tcbpractice}{來動手：看看你的手機有沒有記憶體拓展}
\begin{enumerate}
  \item 開啟【設定】，搜尋"記憶體擴充"、"記憶體融合"、"RAM Plus"或"Extended RAM"。
  \item 找到後檢視當前是否開啟，以及拓展容量設定是多少。
  \item 對照本節內容判斷：你的機型實體記憶體多大？快閃記憶體規格是什麼？剩餘儲存空間是否充足？
  \item 嘗試關閉或調低拓展容量，觀察一週後臺保活和發熱變化，找到最適合自己的設定。
\end{enumerate}

\end{tcbpractice}
